\input{./._relpath-to-latexroot.ltx}
\documentclass[\latexroot/\projectname]{subfiles}
\input{\latexroot/@resources/subfile-setup.ltx}
\begin{document}

\whenintegrated{\label{microeconomic-literature}}\par\subsection{Related literature}
\whenintegrated{\label{sec:microlit}} % integrated doc: SST for labels

This paper relates to the growing body of work that uses HFI to study the macroeconomic effects of monetary policy. Building on early contributions such as \hyperlink{Kuttner (2001)}{Kuttner (2001)} and \hyperlink{Gürkaynak et al. (2005)}{Gürkaynak et al. (2005)}, a large literature combines interest rate surprises with VARs  to estimate the effects of conventional and unconventional MP shocks (e.g., \hyperlink{Gertler and Karadi (2015)}{Gertler and Karadi (2015)}; \hyperlink{Nakamura and Steinsson (2018)}{Nakamura and Steinsson (2018)}; \hyperlink{Andrade and Ferroni (2021)}{Andrade and Ferroni (2021)}; \hyperlink{Miranda-Agrippino and Ricco (2023)}{Miranda-Agrippino and Ricco (2023)}; \hyperlink{Bauer and Swanson (2023a, 2023b)}{Bauer and Swanson (2023a, 2023b)}; \hyperlink{Jarociński and Karadi (2020, 2025)}{Jarociński and Karadi (2020, 2025)}).

Our analysis complements \hyperlink{Jarociński and Karadi (2025)}{Jarociński and Karadi (2025)}, who jointly identify and measure the effects of CBI and RN  components for the U.S., by extending their framework to an EME and by explicitly distinguishing between conventional and unconventional MP shocks. Beyond isolating these components for two types of interest rate surprises, we trace their dynamic effects on a broader set of variables, including long-term bond yields, the exchange rate, and inflation expectations. Notably, we show that RN shocks appear to play distinct roles in the transmission of target-rate and FG shocks.

We also contribute to the literature analyzing the macroeconomic effects of unconventional monetary policy shocks. \hyperlink{Campbell et al. (2012)}{Campbell et al. (2012)} and \hyperlink{Andrade and Ferroni (2021)}{Andrade and Ferroni (2021)} study the macroeconomic effects of forward guidance (FG) in advanced economies, emphasizing that FG surprises may combine “Odyssean” policy commitments with “Delphic” signals about future macroeconomic conditions. \hyperlink{Swanson (2024)}{Swanson (2024)} shows that FG surprises may be contaminated by RN effects. For the case of Mexico, \hyperlink{Solís (2023)}{Solís (2023)} documents that FG significantly affects the exchange rate and the yield curve, underscoring the importance of forward-looking communication channels in a small open economy.

We add to this literature by providing new evidence on the macroeconomic effects of FG in an EME. We find that unpurged FG surprises generate disproportionately large and faster responses—most notably in real output—resembling what has been described as the FG puzzle. Once RN components are isolated, these distortions are substantially attenuated, and the estimated effects of FG become more comparable to those of conventional MP shocks.

Finally, this paper contributes to the literature on monetary transmission in EMEs, where empirical work using HFI has documented  puzzling responses of exchange rates, prices, and inflation expectations to contractionary MP shocks. \hyperlink{Aruoba et al. (2021)}{Aruoba et al. (2021)} show that the domestic currency appreciates and inflation expectations rise following an unexpected tightening in Chile. Using a panel of EMEs, \hyperlink{Bolhuis et al. (2024)}{Bolhuis et al. (2024)} document an average depreciation of the exchange rate after contractionary MP shocks. More recently, \hyperlink{Witheridge (2024)}{Witheridge (2024)} reports increases in inflation after an unexpected tightening and interprets them as evidence of fiscal-led policy dynamics in some EMEs, including Mexico. In contrast, \hyperlink{Checo et al. (2024)}{Checo et al. (2024)} find more conventional transmission patterns using survey-based forecast errors, while \hyperlink{Gomes et al. (2024)}{Gomes et al. (2024)} document sizable real and financial effects of high-frequency tightening shocks in Brazil. For Mexico, \hyperlink{Solís (2023)}{Solís (2023)} shows that exchange-rate puzzles largely disappear when monetary policy surprises are identified using intraday data.

We contribute to this literature by proposing an alternative and complementary explanation for the emergence of these puzzles: the failure to account for informational components embedded in monetary policy surprises. We show that once CBI and RN components are isolated, impulse responses become broadly consistent with standard transmission mechanisms. Importantly, we develop the analysis using daily data, which is particularly valuable in EMEs where access to intraday financial data may be limited.

% Smart bibliography: Only include bibliography if standalone AND has citations
\smartbib

\pdfonly{
% execute this version if compiling with pdflatex
\end{document}
}
